\documentstyle[%


righttag%


,11pt%


,amssymb%


,russian%


,izvru%


]{amsart}





\newcommand{\tl}[3]{\medskip\noindent{\bf#1}\ #2 \dotfill #3 \par} 


\pagestyle{empty}


% \pagestyle{plain}


\begin{document}


\par


\vskip 2cm


\par


\bigskip





В данном номере представлены статьи участников международной


конференции ``Семинар имени Н.И.\,Лобачевского'', которая


проходила с 28 ноября по 1 декабря 2002~года в связи с очередным


присуждением медали им.~Н.И.\,Лобачевского за выдающиеся достижения


в геометрии.











\bigskip


\bigskip


\par


\centerline{\Large СОДЕРЖАНИЕ}


\bigskip


\bigskip


%%%%%%%%%% Use \newline %%%%%%%%%%%








\tl{Акивис М.А., Гольдберг В.В.}


{Многообразия с вырожденным гауссовым отображением\newline с кратными


фокусами и скрученные конусы}{3}


\tl{Джорич Мирьяна, Окумура Масафуми.}


{О кривизне $CR$-подмногообразий комплекс-\newline ного проективного


пространства, имеющих максимальную $CR$-размерность}{15}


\tl{Игошин В.А., Китаева Е.К., Коткова Н.В.}


{Об аффинных симметриях квазигеодезичес-\newline ких потоков}{24}


\tl{Киосак В.А., Микеш Й.}


{О геодезических отображениях пространств Эйнштейна}{36}


\tl{Малахальцев М.А.}


{Инфинитезимальные деформации симплектической


структуры с\newline особенностями}{42}


\tl{Салимов А.А., Ченгиз Н.}


{Поднятие римановых метрик на тензорные расслоения}{51}


\tl{Сидорякина В.В.}


{Бесконечно малые ARG-деформации поверхностей при условии обобщен-\newline ного 


скольжения}{60}


\tl{Степанов С.Е.}


{О голоморфном отображении


почти семи-келерова многообразия}{67}


\tl{Столяров А.В.} 


{Пространство проективно-метрической связности}{70}


\tl{Султанов А.Я.} 


{Аффинные преобразования многообразий с  


линейной связностью и\newline автоморфизмы линейных алгебр}{77}


\tl{Толстихина Г.А., Шелехов А.М.}


{Многоточечные инварианты групп преобразований\newline и определяемые ими


три-ткани}{82}


\tl{Шурыгин В.В.}


{О строении полных многообразий над алгебрами Вейля}{88}











%\bigskip


%\par


%\centerline{\Large КРАТКИЕ~СООБЩЕНИЯ}


%\bigskip











\vfill


\noindent\raisebox{2pt}{\copyright} Казанский госудаpственный унивеpситет


\end{document}


